\documentclass{article}
\usepackage{../math_template}

\author{Jonathan Kasongo}
\title{Canadian Math Olympiad 1971 --- P10}

\begin{document}
\maketitle

\begin{problem}{Canadian Math Olympiad 1971 --- P10}
Suppose that $n$ people each know exactly one piece of information,
and all $n$ pieces are different. Every time person $A$ phones person
$B$, $A$ tells $B$ everything that $A$ knows, while $B$ tells $A$ nothing.
What is the minimum number of phone calls between pairs of people
needed for everyone to know everything? Prove your answer is a minimum.
\end{problem}

\begin{solution}{Solution}
The answer is $2(n-1)$. Finding a construction is easy.
There are two main steps. \vspace{0.2cm}

Firstly we show that it takes at least $n-1$ phone calls to have
1 person know all bits of information. For someone to know
all pieces of information, each of the other $n-1$ people must've
phoned someone else at least one time, otherwise a piece of information
would still be known uniquely by exactly one person. \vspace{0.2cm}

Now once we have someone that knows all pieces of information, for the
remaining $n-1$ people to know everything it takes at least one phone
call for them to know everything. Thus in each case we add up to get
at least $2(n-1)$ phone calls. $\blacksquare$
\end{solution}

\end{document}
